%----------------------------------------------------------------------------------------
%	DOCUMENT DEFINITION
%----------------------------------------------------------------------------------------

\documentclass[a4paper,12pt]{article}

%----------------------------------------------------------------------------------------
%	PACKAGES
%----------------------------------------------------------------------------------------
\usepackage{url}
\usepackage{parskip}

%other packages for formatting
\RequirePackage{color}
\RequirePackage{graphicx}
\usepackage[usenames,dvipsnames]{xcolor}
\usepackage[scale=0.9]{geometry}

%tabularx environment
\usepackage{tabularx}

%for lists within experience section
\usepackage{enumitem}

% centered version of 'X' col. type
\newcolumntype{C}{>{\centering\arraybackslash}X}

%to prevent spillover of tabular into next pages
\usepackage{supertabular}
\usepackage{tabularx}
\newlength{\fullcollw}
\setlength{\fullcollw}{0.47\textwidth}

%custom \section
\usepackage{titlesec}
\usepackage{multicol}
\usepackage{multirow}

%CV Sections inspired by:
%http://stefano.italians.nl/archives/26
\titleformat{\section}{\Large\scshape\raggedright}{}{0em}{}[\titlerule]
\titlespacing{\section}{0pt}{10pt}{10pt}

%Setup hyperref package, and colours for links
\usepackage[unicode, draft=false]{hyperref}
\definecolor{linkcolour}{rgb}{0,0.2,0.6}
\hypersetup{colorlinks,breaklinks,urlcolor=linkcolour,linkcolor=linkcolour}

%for social icons
\usepackage{fontawesome5}

% job listing environments
\newenvironment{jobshort}[2]
    {
    \begin{tabularx}{\linewidth}{@{}l X r@{}}
    \textbf{#1} & \hfill &  #2 \\[3.75pt]
    \end{tabularx}
    }
    {
    }

\newenvironment{joblong}[2]
    {
    \begin{tabularx}{\linewidth}{@{}l X r@{}}
    \textbf{#1} & \hfill &  #2 \\[3.75pt]
    \end{tabularx}
    \begin{minipage}[t]{\linewidth}
    \begin{itemize}[nosep,after=\strut, leftmargin=1em, itemsep=3pt,label=--]
    }
    {
    \end{itemize}
    \end{minipage}
    }

%----------------------------------------------------------------------------------------
%	BEGIN DOCUMENT
%----------------------------------------------------------------------------------------
\begin{document}

% non-numbered pages
\pagestyle{empty}

%----------------------------------------------------------------------------------------
%	TITLE
%----------------------------------------------------------------------------------------

\begin{tabularx}{\linewidth}{@{}l X r@{}}
\Huge{Doyeon Hwang} \\[7.5pt]
\normalsize
\href{https://github.com/fwangdo}{\raisebox{-0.05\height}\faGithub\ fwangdo} \ $|$ \
\href{mailto:fwangdo35@gmail.com}{\raisebox{-0.05\height}\faEnvelope \ fwangdo35@gmail.com} \ $|$ \
\href{tel:+821020121394}{\raisebox{-0.05\height}\faMobile \ +82-10-2012-1394}
\end{tabularx}

\vspace{10pt}

%----------------------------------------------------------------------------------------
% SUMMARY
%----------------------------------------------------------------------------------------
\section{Summary}
Compiler engineer with hands-on experience building a deep learning compiler from frontend IR design through MLIR-based middle-end lowering.
Strong interest in applying program analysis techniques to optimize computational graphs and lower them to efficient executable forms.
Background in program testing and verification, with a track record of finding 40+ bugs in production SMT solvers.

%----------------------------------------------------------------------------------------
%	EDUCATION
%----------------------------------------------------------------------------------------
\section{Education}

\begin{tabularx}{\linewidth}{@{}l X@{}}
2022 -- 2025 & \textbf{Ph.D. Coursework Completed} (Computer Science), Korea University \hfill Seoul, Korea \\
& \small{SMT Solver Fuzzing, Abstract Interpretation} \\[3pt]

2019 -- 2021 & \textbf{M.S.} (Information Convergence Engineering), Pusan National University \hfill Busan, Korea \\[3pt]

2012 -- 2018 & \textbf{B.A.} (Psychology), Jeonbuk National University \hfill Jeonju, Korea \\
\end{tabularx}

%----------------------------------------------------------------------------------------
%	PROJECTS
%----------------------------------------------------------------------------------------
\section{Projects}

\begin{tabularx}{\linewidth}{ @{}l r@{} }
\textbf{PyTorch-based Deep Learning Compiler Frontend}
& \hfill \href{https://github.com/fwangdo/Gawee}{GitHub} \\[3.75pt]
\multicolumn{2}{@{}X@{}}{
\begin{minipage}[t]{\linewidth}
    \begin{itemize}[nosep,after=\strut, leftmargin=1em, itemsep=2pt,label=--]
        \item Designed a custom IR and converted PyTorch models to it via FX graph tracing
        \item Implemented graph-level rewrite passes (Conv--BatchNorm folding, Constant folding, etc.)
        \item Quantitatively verified optimization effects by comparing memory access cost estimates before/after optimization (ResNet-18, UNet)
    \end{itemize}
    \end{minipage}
}
\end{tabularx}

\vspace{6pt}

\begin{tabularx}{\linewidth}{ @{}l r@{} }
\textbf{MLIR-based Deep Learning Compiler Middle-end (C++)}
& \hfill \href{https://github.com/fwangdo/Gawee}{GitHub} \\[3.75pt]
\multicolumn{2}{@{}X@{}}{
\begin{minipage}[t]{\linewidth}
    \begin{itemize}[nosep,after=\strut, leftmargin=1em, itemsep=2pt,label=--]
        \item Defined a custom MLIR Dialect and operations (Conv, ReLU, Add, etc.) using TableGen
        \item Built an Emitter that converts frontend JSON IR into MLIR operations
        \item Wrote OpConversionPattern-based lowering patterns from custom Dialect to Linalg
    \end{itemize}
    \end{minipage}
}
\end{tabularx}

\vspace{6pt}

%----------------------------------------------------------------------------------------
%	RESEARCH EXPERIENCE
%----------------------------------------------------------------------------------------
\section{Research Experience}

\begin{joblong}{SMT Solver Testing Tool Development (Python)}{2022 -- 2025}
\item Detected \textbf{40+ bugs} in state-of-the-art SMT solvers (Z3, CVC5, bitwuzla) using a custom fuzzer
\item Designed a custom IR to unify I/O formats across heterogeneous solvers; implemented an SMTLibv2 parser using PyParsing
\item Engineered a fuzzer that generates syntactically valid test cases while covering diverse production rules
\end{joblong}

\vspace{6pt}

%----------------------------------------------------------------------------------------
%	SKILLS
%----------------------------------------------------------------------------------------
\section{Skills}

\begin{tabularx}{\linewidth}{@{}l X@{}}
Languages & \normalsize{Python, C++, JavaScript, OCaml} \\
Tools \& Technologies & \normalsize{Git, Linux, MLIR, PyTorch, PyParsing, LaTeX} \\
Domains & \normalsize{Compilers (Frontend / Middle-end), Program Analysis, Program Testing} \\
\end{tabularx}

\vfill

\end{document}
